%----------------------------------------------------------------------------------------------
\chapter{Introducción a git \& github}
%----------------------------------------------------------------------------------------------
\minitoc% Creating an actual minitoc

\begin{justify}
	\begin{multicols}{2}		
		{\bf ¿Qué es Git?}
		
		Git es una herramienta de código abierto que permite a los desarrolladores llevar
		un registro de los cambios que hacen en su código a lo largo del tiempo. Es como
		una máquina del tiempo para tu código, ya que puedes ver quién hizo qué, cuándo
		lo hizo y porqué. Esto es crucial para la colaboración en equipo, ya que permite
		a varias personas trabajar en el mismo proyecto sin sobrescribir los cambios de los
		demás. Con Git, los desarrolladores trabajan en su propia copia local del proyecto
		(repositorio), y luego sincronizan sus cambios con el repositorio central.\\
		
		{\bf ¿Qué es GitHub?}
		
		Github es una plataforma de hosting para repositorios de Git. Piénsalo como una red social
		para desarrolladores. Permite a los equipos almacer su código en un sólo lugar centralzado
		en la nube, lo que facilita la colaboración y el intercambio de código. GitHub también 
		ofrece funciones adicionales como seguimiento de errores (Issues), gestión de proyectos
		(Projects), y herramientas de integración continua/despliegue continuo (CI/CD) a través 
		de GitHub Actions. 
		Además, la naturaleza pública de muchos repositorios en GitHub promueve el desarrollo de 
		software de código abierto.
	\end{multicols}
\end{justify}


{\bf Tres ejemplos de uso}

\begin{justify}
	\begin{enumerate}
		\item {\bf Desarrollo de una aplicación web en equipo}: Un equipo de 10 desarrolladores está creando un sitio
		web. Uno trabaja en la página de inicio, otro en la base de datos, y un tercero en el diseño de la 
		interface. Usando Git, cada desarrollador trabaja en su propia rama local de código. Luego, suben sus
		cambios a un repositorio central de GitHub. GitHub gestiona los {\bf pull requests} (solicitud de extracción)
		para revisar el código, lo que asegura que todos los cambios se integren sin conflictos, manteniendo
		un historial claro de quién contribuyó con qué.
		\item {\bf Manejo de versiones de un documento de investigación}: Un investigador está escribiendo una tesis
		y quiere guardar un historial completo de cada cambio que hace. Usando Git localmente, puede crear {\bf commits}
		(instantáneas) en diferentes puntos del proceso de e scritura, etiquetando las versiones importantes
		(por ejemplo, Versión 1.0- borrador inicial, Versión 2.5 - revisada por el supervisor). Si comete un error, 
		puede volver a una versión anterior sin perder trabajo. Si quisiera colaborar con otros investigadores, podría
		subir el proyecto a GitHub para que ellos también puedan contribuir y revisar el documento.
		\item {\bf Contribución a un proyecto de código abierto}: Un desarrollador encuentra un error en una biblioteca
		de software popular y quiere solucionarlo. El código fuente de la biblioteca está alojado en GitHub. El
		desarrollador {\bf forkea} (hace una copia) del repositorio a su propia cuentra de GitHub, realiza los 
		cambios necesarios en su copia local usando Git, y luego envía un {\bf pull request} al repositorio original.
		Los mantenedores del proyecto revisan su código y, si es aceptado, lo fisionan con la base de código principal,
		beneficiando a todos los usuarios de la biblioteca.
	\end{enumerate}
\end{justify}


\textcolor{red}{\bf Importante!}

El contenido de este capítulo es una una copia ó adaptación de las siguientes fuentes:

\begin{itemize}
	\item[$\bullet$] \href{https://yenchiah.me/jupyter-book-template/docs/home.html}{Jupyter Book Tutorial}
	\item[$\bullet$] \href{https://jupyterbook.org/en/stable/start/build.html}{Build your book}
\end{itemize}

%----------------------------------------------------------------------------------------------
\section{Instalar Git}
%##############################################################################################
\subsection{Arch Linux.}

\commandbox*{sudo pacman -Syu}

Instalar Git.

\commandbox*{sudo pacman -S git}

Verificar la instalación.

\commandbox*{git --version}

\subsection{Mac.}

\subsection{Windows.}


%----------------------------------------------------------------------------------------------
\section{Configuración inicial}
%----------------------------------------------------------------------------------------------
\subsection{Configurar nombre de usuario}
\commandbox*{git config --global user.name "Nombre de Usuario"}

\subsection{Configurar correo electrónico}
\commandbox*{git config --global user.email "tu.email@ejemplo.com"}

\subsection{Configurar editor de texto}
\commandbox*{git config --global user.editor "nvim"}

\subsection{Verificar configuración}
\commandbox*{git config --list}


\begin{center}
	\begin{tcolorbox}[space to upper,
		enhanced, 
		%colback=green!5!white,colframe=green!75!black,
		%colback=purple!5!white,colframe=purple!75!black,
		%colback=black!5!white,colframe=black!75,
		colback=orange!5!white,colframe=orange!95,
		%skin=bicolor,	
		boxrule=0.15pt,
		collower=white,
		%title={01\_hello-rust.rs},
		halign=center,
		valign=center,
		nobeforeafter,
		halign lower=flush right,
		top=1mm,left=1mm,right=1mm, bottom=0mm,
		%lifted shadow={1mm}{-2mm}{3mm}{0.1mm}{black!50!white},
		text width=6cm,
		height=7cm
		]
	\begin{forest}
		pic dir tree,
		pic root,
		for tree={% folder icons by default; override using file for file icons
			directory,
			fit=band,
		},
		[hello-rust
		[target
		[doc]
		[debug
		[hello-rust, file]
		]
		]
		[src
		[main.rs, file]
		]
		[Cargo.lock,file]
		[Cargo.tolm,file]
		]
	\end{forest}
	\end{tcolorbox}
\end{center}

Finalmente, la instrucción \commandbox*{./hola-mundo}, pone en marcha el ejecutable \commandexe{hola-mundo}

Por ejemplo, el código \commandfilename{hola-mundo.rs}, es una aplicación básica que muestra un mensaje. 


Para compilar el código que se encuentra en el archivo \commandfilename{hola-mundo.rs}, se ejecuta la instrucción \commandbox*{rustc hola-mundo.rs}, el cuál generá un archivo ejecutable \commandexe{hola-mundo}.

\section{ssh-keygen}
\commandfilename{ssh-keygen} es una utilidad que forma parte del paquete{\bf OpenSSH} en Linux. Su función
principal es {\bf generar y administrar pares de claves} (públicas y privadas) para el protocolo {\bf SSH}
(Secure Shell). Estas claves se utilizan para la autenticación sin contraseña entre clientes y servidores, 
mejorando la seguridad y automatización.

\section{Trabajar un mismo repositorio desde diferentes computadoras}
%##############################################################################################

\subsection{Clonar el repositorio en cada computadora}

\commandbox*{git clone https:github.com/nombre_usuario/nombre_repositorio.git}

Por ejemplo:

\commandbox*{git clone https://github.com/vrrp/indicesBioclimaticos.git}

Donde, nombre\_usuario: vrrp, y nombre\_repositorio: indicesBioclimaticos
